% !TEX root = ../discretemath.tex

\section{Вершинная связность}
\begin{Definition}{Разделяющее множество}{}
    В графе $G$ множество вершин $R\subset V$ называется разделяющим для множеств 
    $X\subset V,Y\subset V$($X$ и $Y$ могут пересекаться), если в графе $G$, индуцированном
    на $V\setminus R$, нет путей между вершинами из $X\setminus R$ и $Y\setminus R$.
\end{Definition}
\begin{Theorem}{Геринг}{}
    Дан граф $G$ и два множества его вершин $X,Y$. 
    Если любое множество, разделяющее $X$ и $Y$, состоит хотя бы из $k$ вершин,
    то существует $k$ непересекающихся по вершинам $XY$-путей
    (от вершины множества $X$, до вершины множества $Y$).
\end{Theorem}
\begin{Statement}{}{}
    Корректна следующая перефорулировка теоремы:
    наибольшее количество непересекающихся по вершинам $XY$-путей
    равно минимальному размеру множества, разделяющего $X$ и $Y$.
\end{Statement}
\begin{proof}
    Будем доказывать наличие $k$ непересекающихся по вершинам $XY$-путей в графе,
    где все $XY$-разделяющие множества содержат хотя бы $k$ вершин.

    Доказательство по индукции по количеству вершин.

    База: $\left|V\right|=k$. Очевидно, что в таком случае $V\setminus R=\emptyset$.

    Переход:

    Случай первый. $\exists R$~--- $XY$-разделяющее множество такое, что
    $\left|R\right|=k, X\setminus R\neq\emptyset,Y\setminus R\neq\emptyset$.
    Нетрудно понять, что $X\neq Y$.
    $V_X$~--- вершины, до которых можно дойти из $X\setminus R$ в графе $G\setminus R$.
    Аналогично определим $V_Y$. $G_X=G\mid_{V_X\cup R},G_Y=G\mid_{V_Y\cup R}$
    \begin{Statement}{}{}
        Любое множество, разделяющее $X$ и $R$ в $G_X$, также является разделяющим множеством
        для множеств $X$ и $Y$ в $G$.
    \end{Statement}
    \begin{proof}
        Пусть не так, тогда зафиксируем $XR$-разделяющее множество $S$.
        При этом существует путь $p$ из $X$ в $Y$, не содержащий вершин из $S$.
        Так как $R$ разделяет $X$ и $Y$, то путь $p$ посетит хотя бы одну вершину из $R$.
        Следовательно, $S$ не является $XR$-разделяющим множеством.
        
        Замечание: $R$ разделяет в $G$, $S$ разделяет в $G_X$,
        а путь $p$ проходит по $G\setminus S$. Так что в теории он может пролегать
        по нескольким вершинам из $R$ и несколько раз переходить из множества $V_X$ в $V_Y$ и обратно.
        Придется уточнить, что именно первая встретившаяся по пути $p$ вершина из $R$ порождает противоречие.
    \end{proof}
    Из вышедоказанного получаем, что любое множество разделяющее $X$ и $R$ в $G_X$ содержит хотя бы $k$ вершин.
    Из индукционного предположения знаем, что в $G_X$ есть $k$ непересекающихся $XR$-путей.
    Аналогично в $G_Y$ есть $k$ непересекающихся $YR$-путей.
    Объединим $XR$-пути и $YR$-пути, получим $k$ непересекающихся $XY$-путей.

    Случай второй. $\nexists$ вышеописанного $R$

    Случай второй, подслучай первый. $X\cup Y\neq V\implies\exists u\not\in V\setminus(X\cup Y)$
    \begin{Statement}{}{}
        В $G\setminus\left\{u\right\}$ нет $XY$-разделяющего множества $S$ размера $k-1$, такого что 
        $X\setminus S\neq\emptyset,Y\setminus S\neq\emptyset$
    \end{Statement}
    \begin{proof}
        Пусть такое множество есть, тогда добавим в него $u$,
        получим $XY$-разделяющее множество размера $k$.
        Такого множества не может быть при рассмотрении случая 2.
    \end{proof}
    Если в $G\setminus\left\{v\right\}$ нет $XY$-разделяющего множества размера $k-1$, значит там все множества
    хотя бы размера $k$. По индукционному предположению в $G\setminus\left\{v\right\}$ есть $k$ путей, значит и в $G$ есть $k$ путей.

    Случай второй, подслучай второй. $X\cup Y = V$ Очевидно, что $X\cap Y$ будет входить в любое разделяющее множество.
    $G\setminus\left(X\cap Y\right)$ можно представить как двудольный граф, где в левой половине вершины из $X$, а в правой из $Y$.
    Возьмем минимальное вершинное покрытие этого графа и добавим его в разделяющее множество. Получим минимальное разделяющее множество.
    Получаем $k=\left|\text{минимальное разделяющее}\right|=\left|\text{минимальное вершинное покрытие}\right|+\left|X\cap Y\right|
    =\left|\text{максимальное паросочетание}\right|+\left|X\cap Y\right|=\left|\text{пути }X\to Y\right|$.
\end{proof}
\begin{Theorem}{Менгер}{}
    Количество непересекающихся по внутренним вершинам путей между несвязными вершинами $v$ и $u$ равно размеру минимального 
    $vu$-разделяющего множества $S:S\subset V\setminus\left\{v,u\right\}$.
\end{Theorem}
\begin{proof}
    Построим граф $G'$. В нем заменим вершину $v$ её клонами в количестве $k$ штук. Аналогично заменим $u$.
    Теперь по теореме Геринга количество путей между множествами клонов $v$ и $u$ равно размеру минимального разделяющего множества.
\end{proof}
\begin{Definition}{}{}
    Граф $G$ называется $k$-связным, если $\left|V(G)\right|>k$ и $\forall S\subset V(G),\left|S\right|<k,G\setminus S$ связен.
\end{Definition}
\begin{Definition}{}{}
    $\kappa_G(v,u)$ ~---количество непересекающихся по внутренним вершинам путей от $v$ до $u$ в графе $G$.
\end{Definition}
\begin{Theorem}{}{}
    $G$ $k$-связен $\iff\forall v\neq u:\kappa_G(v,u)\geq k$
\end{Theorem}
\begin{proof}
    $\Rightarrow$ Для несвязных $v,y$ по теореме Менгера. Для связных вершин добавим в граф в $G$ две вершины $v'$ и $u'$.
    $v'$ соединим с соседями $u$, а $u'$ с соседями $v$. Связность графа не понизится, по Менгеру есть $k$ непересекающихся путей
    от $v'$ до $u'$. Теперь преобразуем эти пути так, чтобы они были в изначальном графе и между вершинами $v$ и $u$:
    \begin{itemize}
        \item[]если в пути есть вершины $v$ и $u$, вместо этого пути возьмем путь $\left\{v,u\right\}$;
        \item[]если в пути есть только вершина $v$, вместо этого пути возьмем его отрезок от $u'$ до $v$ и заменим в нем $u'$ на $u$;
        \item[]если в пути есть только вершина $u$, вместо этого пути возьмем его отрезок от $u$ до $v'$ и заменим в нем $v'$ на $v$;
        \item[]если в пути нет вершин $u$ и $v$, то заменим в нем $v'$ на $v$, а $u'$ на $u$.
    \end{itemize}

    $\Leftarrow$ пусть не так, тогда $\exists S,\left|S\right|<k,G\setminus S$ не связен.
    Тогда $\exists a,b\in G\setminus S$, которые лежат в разных компонентах связности.
    Вершины $a$ и $b$ несмежные и $\kappa_G(a,b)<k$, так как удаление $S$ нарушило их связность.
    Противоречие.
\end{proof}
\begin{Definition}{}{}
    $v\in V(G)$~--- точка сочленения, если $G\setminus\left\{v\right\}$ несвязен.
\end{Definition}
\begin{Definition}{}{}
    $B$~---блок, если:
    \begin{itemize}
        \item[]$B$~---индуцированный подграф на какое-то множество вершин
        \item[]$B$~---ребро или двусвязный граф
        \item[]$\nexists B'\supset B:B'$ удовлетворяет первым двум признакам
    \end{itemize}
\end{Definition}
$A$~---множество точек сочленения, $B$~---множество блоков. Строим граф из $A\cup B$. Проводим ребра между $a\in A и b\in B$ тогда, когда $a\in b$.
\begin{Theorem}{}{}
    Полученный граф является деревом.
\end{Theorem}
\begin{proof}
    \begin{itemize}
        \item[]У двух разных блоков не может быть более одной общей вершины, иначе их объединение является блоком.
        \item[]Аналогично в графе не может быть цикла из блоков, иначе их объединение является блоком.
    \end{itemize}
    
\end{proof}
